\documentclass[11pt]{article}
\usepackage[T1]{fontenc}
\usepackage{lmodern,amsmath,amssymb,booktabs}
\usepackage[a4paper,margin=23mm]{geometry}
\usepackage[hidelinks]{hyperref}
\hypersetup{pdftitle={Unbounded gaps between the two largest distinct Smith factors of a one-defect connectivity matrix},pdfauthor={Iaroslav Argunov}}
\newtheorem{theorem}{Theorem}
\newtheorem{lemma}{Lemma}
\newenvironment{proof}{\par\noindent\textit{Proof.}\ }{\hfill\(\square\)\par}
\newcommand{\F}{\mathbb F_2}
\newcommand{\val}{v_2}
\newcommand{\ord}{\operatorname{ord}_{\pi}}
\newcommand{\im}{\operatorname{im}}
\newcommand{\Tr}{\mathsf{T}}
\begin{document}
\begin{center}
{\Large\bfseries Unbounded gaps between the two largest distinct Smith factors of a one-defect connectivity matrix\par}
\medskip Iaroslav Argunov\\[2pt] 19 November 2026
\end{center}
\begin{abstract}
Sum connected gluings of marked noncrossing matchings over a subgroup
of rotations on \(2k+1\) sites. For \(k\ge3\) the quotient of the two
largest distinct integer Smith factors is the largest power of two
dividing the even member of \(\{k,k+1\}\). We determine both multiplicities
for every subgroup. The gap follows from complementing the incidence
generators of the dual lattice. The surviving form at its far end is a
binary incidence pairing involving only points, pairs, triples and four-sets.
\end{abstract}

\section{The gap theorem}
Label \(N=2k+1\) points on an oriented circle. A state is a marked site and
a noncrossing perfect matching of the others. Let \(A(p,q)\) be the
indicator that their union is connected, retaining both copies of any
common edge. For a rotation subgroup \(H\) of order \(h\), choose one
state per orbit and form \(G_H(p,q)=\sum_{g\in H}A(p,gq)\).
Set \(m=N/h,\ e=\val((k+1)!)\). Write \(t\) for the number of
\(H\)-orbits of three-sets. Burnside's lemma gives
\[
 t=\bigl(\tbinom N3+{\bf1}_{3\mid h}\,2N/3\bigr)/h.
\]
Only the two elements of order three can fix a triple besides the identity.

\begin{theorem}\label{thm}
For \(k\ge3\) put \(a=\val(k(k+1))\). The two largest distinct Smith
factors of \(G_H\) are \(2^e\) and \(2^{e-a}\). Their multiplicities are
\[
\begin{array}{c|cc}
 &2^e&2^{e-a}\\ \hline
k\ {\rm odd}&1&m(k+1)-2\\
k\ {\rm even}&m+1&t+m(k-2)-2.
\end{array}
\]
For \(k=0\) the only factor is one. At \(k=1,2\) the largest
factor is two, with respective multiplicities \(1,m+1\), and the
remaining \(m-1\) factors are one in either case.
\end{theorem}
For full rotations, \(k=31\) gives a single \(2^{31}\) followed by
30 copies of \(2^{26}\). At \(k=32\) there are two copies of \(2^{31}\)
followed by 700 copies of \(2^{26}\). In either case all four intermediate
powers are absent. Orders \(k=2^b-1\) or \(2^b\) give gaps of unbounded
length in the sequence of exponents.

\section{From marked diagrams to an integral form}
Take \(R=\mathbb Z[i]_{(1+i)}\), let \(\pi=1+i\), and normalize
\(\ord(2)=2\). In the permutational coordinate space on the \(k\)-subsets,
identify the vector of a subset with its square-free monomial; write
\(r_S\) for its containment test. On a marked diagram \(p\), orient
each arc \(a\to b\) along the boundary interval avoiding the mark and put
\[
 T(p)=\prod_{a\to b}(x_b-i x_a).
\]
There are \(N\binom{2k}k/(k+1)=\binom Nk\) diagrams, so \(T\) is square.
In Definition 3.1, Eq.~(12), of Morin-Duchesne and Saint-Aubin, set
\(v=1\) and \(u=e^{\pi i/4}\). Each arc contributes
\(u x_b+u^{-1}x_a=u(x_b-i x_a)\), so the published map equals \(u^kT\).
In the one-defect sector the open path joins the two ends. An essential
closed component would separate the ends and intersect that path.
Thus all closed components are
contractible. At twist \(v=1\) and loop weight \(\beta=0\),
Theorem 3.3 and Proposition 4.3, Eq.~(29), of
\href{https://arxiv.org/abs/1203.4996}{Morin-Duchesne and Saint-Aubin}
give
\[
 A=i^k T^{\mathrm T}T,\qquad |\det A|=2^E,\qquad
 E=\sum_{j=0}^{k-1}\binom Nj .
\]
In the determinant each factor is
\(2^2-4\cos^2((r+1/2)\pi/2)=2\); the transpose is bilinear.
The mark makes the action of \(H\) on diagrams free. The orbit sums
are an integral basis of a primitive invariant sublattice
of the free diagram module. The operator \(A\) acts there by \(G_H\);
extending its basis makes \(A\) block triangular. Hence \(\det G_H\)
divides \(\det A\) over the integers, and has no odd prime divisor.

We verify the integral image needed below. Define
\(L=\{x:\langle x,r_S\rangle\in\pi^{k-|S|}R,\ |S|\le k\}\)
inside the subset lattice. In \(T(p)\) the test \(r_S\) either vanishes
or selects one endpoint from each of \(|S|\) arcs; it obtains a factor
\((1-i)^{k-|S|}\) from the free choices. Thus \(\im T\subset L\).
The containment rows on the good sets \(s_1<\cdots<s_j\) with
\(s_i\ge2i\) form the integral subset basis of
\href{https://doi.org/10.1006/eujc.1993.1001}{Bier}.
For clarity, split the columns on membership of the last site.
Rows containing it reduce to the good-row basis on \((k-1)\)-sets.
Rows omitting it restrict to the basis on \(k\)-sets of the smaller
ground set if \(N>2k\). At \(N=2k\) complement those \(k\)-sets;
the rows become avoidance tests on \((k-1)\)-sets, related to containment
by triangular inclusion-exclusion since subsets of a good set are good.
This proves the basis assertion inductively from \(k=0\).
There are \(\binom Nj-\binom N{j-1}\) good sets of size \(j\) by ballot
reflection. The rows of size at most \(j\) have rational rank at most
\(\binom Nj\), by summing all \(j\)-extensions of any one row. The
good rows up to that degree already span a saturated space of that rank.
Every test of degree \(j\) is therefore an integral combination of good
rows of degrees at most \(j\). In particular, imposing the divisibility
conditions on the good rows imposes them on all the containment tests.
Consequently in dual basis coordinates the colength of \(L\) equals
\(\sum_{j=0}^k(k-j)(\binom Nj-\binom N{j-1})=E\).
This also equals \(\ord(\det T)\) by the Gram identity, so \(L=\im T\).

Since \(T\) is injective and commutes with rotations, the image of the
invariant state lattice is \(L^H\). Its orbit-sum pairing has matrix
\(hG_H\); the factor \(h\) is a unit. Before taking invariants the
dual of \(L\) is spanned by the scaled tests
\(\pi^{s-k}r_S,\ s=|S|\le k\), by the defining congruences and
biduality (the size-\(k\) tests include the coordinate basis).
The projection \(P=h^{-1}\sum_{g\in H}g\) splits the lattice
orthogonally, so the dual of \(L^H\) is spanned by their projections.
A generating family computes the Gram
invariants on a basis, with extra zero factors, by splitting its free
presentation over the discrete valuation ring.

\section{Complement generators and the first residual layer}

Replace \(r_S\) by the indicator \(b_S\) of \emph{avoiding} \(S\).
Inclusion-exclusion acts triangularly on the scaled generators, with
unit diagonal and coefficients \(\pm\pi^{s-|U|}\) at \(U\subset S\).
Hence the generated lattice is unchanged. Scale its form by \(2^e\).
Up to a common unit the resulting entries are
\[
 K_{ST}=\pi^{s+t'}\,h^{-1}\sum_{g\in H} B_{|S\cup gT|},\qquad
 B_j=\frac{(k+1)_j}{(2k+1)_j},\quad t'=|T|.
\]
Here \((x)_j\) denotes the falling factorial and \(B_0=1\). The common
factor is a unit since \(\val\binom{2k+1}{k}=k-e\).
The nonzero exponents of the scaled presentation, counting exponent
zero as well, are \(2(e-v)\) when the original Smith factor is \(2^v\).
This follows by dualizing a basis, since the form on the dual has the
inverse Gram matrix up to a unit. The remaining argument determines
the first two distinct valuations in \(K\).

\begin{lemma}\label{lem}
For odd \(\ell\ge1\) put \(d=\val(\ell+1)\) and
\(D_j=(\ell+1)_j/(2\ell+1)_j\), \(1\le j\le2\ell+1\).
Then \(\ord(\pi^jD_j)\ge2d\). Equality occurs exactly at \(j=2,4\)
within \(j\le\ell+1\).
\end{lemma}
\begin{proof}
For \(D_j\ne0\) write \(b=\lfloor j/2\rfloor\). The even factors of
the denominator contribute \(2^b(\ell)_b\). Cancelling gives
\[
 j+2\val D_j=2d+j-2b+2\val((\ell-b)_{j-1-b}).
\]
For equality \(j=2b\), and \(b-1\) successive factors must be odd.
This requires \(b=1,2\), and both possibilities do give equality.
Past \(\ell+1\) the numerator vanishes.
\end{proof}

In characteristic two, invariant functions on labelled subsets have
orthonormal orbit-indicator coordinates. Orbit sizes and stabilizers
are odd. Thus a parity sum over relative group elements is the
restriction of the labelled incidence operator in those coordinates.

\paragraph{Odd \(k\ge3\).}
Split off the empty generator, a unit pivot. Its Schur complement has
entries obtained by replacing \(B_u\) by \(B_u-B_sB_{t'}\) in \(K\),
for \(s,t'\ge1\). By Lemma~\ref{lem} with \(\ell=k\) the product term
has valuation at least \(4a\) after scaling. The first term has valuation
at least \(2a\), since \(s+t'\ge u\). At valuation \(2a\) the sets must
be disjoint and \(s+t'=2\) or 4. Averaging preserves this lower bound.
Thus at the first possible residual depth only singleton--singleton,
singleton--triple, and pair--pair blocks remain, each counting disjointness.

\paragraph{Even \(k\ge4\).}
Here \(a=\val(k)\) and \(B_2=(k+1)/(2(2k+1))\) has valuation \(-2\)
in \(R\). For \(u\ge2\), conditioning on two avoided sites gives
\[
 B_u=B_2\frac{(k-1)_{u-2}}{(2k-1)_{u-2}}
     = B_2\sum_{j=0}^{u-2}(-1)^j\binom{u-2}{j}
                    \frac{(k)_j}{(2k-1)_j}.
\]
The second identity is inclusion-exclusion for sampling without
replacement. Lemma~\ref{lem} with \(\ell=k-1\) shows that
\(\pi^{s+t'}(B_u-B_2)\) has valuation at least \(2a\) whenever
\(s+t'\ge u\ge2\): the \(j\)th nonconstant summand has valuation at
least \(2a+s+t'-j-2\), including the factor \(B_2\).
Equality is possible only at
\(s+t'=u=j+2=4,6\). At these indices the last summand has coefficient
one and all preceding summands have strictly larger valuation, so
both equalities hold. The possible sums \(4,6\) remain available at
\(k=4\), even when \(B_6=0\).

Retain the empty set, the sum-zero singleton coordinates, and one
reference pair \(E\). The singleton orbit sum equals
\(\pi(k+1)/h\) times the empty generator, and \(m\) is a unit in \(R\).
Modulo \(\pi\) the singleton form is \(I+J\) (coincident sites give
zero, disjoint sites one after scaling). On the sum-zero subspace this
is a nonsingular pairing. Its links to the empty set and reference pair
vanish modulo \(\pi\): the empty link is constant, and \(s+t'=3\) has
valuation at least one. The empty--pair link is the unit \(\pi^2B_2\);
the pair diagonal is zero modulo \(\pi\), since \(s+t'=4\) contributes
valuation at least two even to the constant term. This gives a unit
block of rank \(m+1\).

For every other size \(s\ge2\) replace its projected generator \(u_S\)
by \(u_S-\pi^{s-2}u_E\). Any pairing involving one of these differences
has its constant \(B_2\) term cancelled exactly. The terms against \(E\)
in the expanded difference have their natural degree sum increased
by \(s-2\) or \(t'-2\), which are nonnegative. Hence every remaining link
to a pivot and every residual pairing is divisible by \(\pi^{2a}\).
More explicitly, a mixed entry against a projected pivot of size \(t'\)
is the difference of the two group averages of
\(\pi^{s+t'}(B_{|S\cup gT|}-B_2)\) and
\(\pi^{s+t'}(B_{|E\cup gT|}-B_2)\). In each average the union has size
at least two and at most the unscaled sum of sizes. The same argument
applies separately to the four terms of a difference--difference entry.
The Schur correction between residual generators is divisible by
\(\pi^{4a}\). After division by \(\pi^{2a}\) and reduction modulo \(\pi\),
the pair coordinates are the sum-zero space on pair orbits.
For generators of size greater than two the reference correction
has positive excess degree and disappears. The only surviving
blocks are disjointness on pair--pair, pair--four-set, and
triple--triple. The first two are restricted on the pair side to
sum-zero coordinates.

\section{Ranks at the far end of the gap}
Here are details of the binary rank calculation for both parities.
Write \(D\) for point--pair incidence, \(p=mk\) for the number of pair
orbits, and \(J\) for an all-ones map of the appropriate size.
Since \(N\) is odd, \(DD^{\mathrm T}=I+J\) on points; the image of \(D\)
is their sum-zero subspace of dimension \(m-1\). Thus \(D^{\mathrm T}D\)
is a projector of rank \(m-1\). These identities hold on the invariant
spaces by the orbit-coordinate observation above. Pair disjointness is
\(I+J+D^{\mathrm T}D\), by expanding the indicator of zero intersection.

For odd \(k\), the all-ones map on pairs is a rank-one projector because
\(p\) is odd; it annihilates \(D^{\mathrm T}D\) in both orders. The
pair block consequently has rank \(p-m\). Against triples a singleton
vector \(w\) gives \(\sum_O w_O+\sum_{x\in T}w_{Hx}\). Comparing triples
differing by a point shows that its kernel consists of the constants,
which indeed give zero since \(m\) and three are odd. The singleton
block \(I+J\) annihilates constants, and the triple diagonal block is
zero. This part has rank \(2(m-1)\): split off the perfect pairing
between complementary row and column spaces of the off-diagonal map.
The total rank is \(m(k+1)-2\).

For even \(k\), set \(Q=I+D^{\mathrm T}D\) and let \(W\) be the
sum-zero space on pairs. Pair disjointness restricts to \(Q\) on \(W\).
The map from \(w\in W\) to four-sets \(T\) sums the weights of edges
disjoint from \(T\). Inclusion-exclusion gives
\(\sum_{x\in T}(Dw)_x+\sum_{\{x,y\}\subset T} w_{xy}
=\sum_{\{x,y\}\subset T}(Qw)_{xy}\), since each vertex occurs three
times in the last edge sum. The only edge functions whose sum vanishes
on all four-sets are constants. Indeed comparing two tests differing
in points \(x,y\) gives zero sum of the values \(f_{xi}+f_{yi}\) over
every triple outside \(x,y\). Such a function is constant by exchanging
one index, and is zero by taking a triple. Thus all edges at each
vertex have equal weights. Conversely six equal edge weights sum to zero.
It follows that the kernel on \(W\) is
\(\im D^{\mathrm T}\oplus\langle{\bf1}\rangle\) of dimension \(m\).
Here \(Q{\bf1}={\bf1}\), \(\ker Q=\im D^{\mathrm T}\), and the
edge constant is sum-zero since \(p\) is even. The pairing \(Q\)
vanishes when both arguments lie in this kernel: the image under \(Q\)
is constant, and all elements of \(W\) have sum zero.
Splitting the off-diagonal pairing therefore gives
rank \(2(p-m-1)\).

Finally write \(E,C\) for pair--triple and point--triple incidence.
Disjointness of triples is \(I+E^{\mathrm T}E+
C^{\mathrm T}(I+J)C\): expand the parity sum of the binomial coefficients
of their intersection size. Counting containments gives
\[
 EE^{\mathrm T}=Q,\quad DE=0,\quad CC^{\mathrm T}=I+J,\quad EC^{\mathrm T}=J .
\]
The third equality uses \(\binom{N-1}2=k(2k-1)\) even. Thus
\(\im E=\ker D\); \(E^{\mathrm T}E\) is a projector of rank \(p-m+1\).
The last projector \(C^{\mathrm T}(I+J)C\) has rank \(m-1\) and
annihilates \(E^{\mathrm T}E\) in both orders. The triple block
has rank \(t-p\), giving \(t+m(k-2)-2\) in all. Both residual ranks
in Theorem~\ref{thm} are positive for \(k\ge3\): in the even case
\(t\ge\binom N3/h\) and \(N\ge9\). The small cases follow by direct
gluing or from the dimension and unit block. This proves the theorem.

\paragraph{Computational checks.}
The accompanying program constructs the geometric matrices for every
rotation subgroup at \(1\le k\le5\), with an additional comparison
against fourteen stored geometric profiles. It checks the phase in
the spin identity on 1,962 gluings. For \(N=63,65\) it computes the
ranks of the residual orbit-incidence forms directly. The proof above
applies at arbitrary order without these computations.
\begin{thebibliography}{9}
\bibitem{bier} T. Bier, Remarks on recent formulas of Wilson and Frankl,
\emph{European J. Combin.} \textbf{14} (1993), 1--8.
\href{https://doi.org/10.1006/eujc.1993.1001}{doi:10.1006/eujc.1993.1001}.
\bibitem{mdsa} A. Morin-Duchesne and Y. Saint-Aubin,
A homomorphism between link and XXZ modules over the periodic Temperley--Lieb algebra,
\emph{J. Phys. A} \textbf{46} (2013), 285207.
\href{https://doi.org/10.1088/1751-8113/46/28/285207}{doi:10.1088/1751-8113/46/28/285207}.
\end{thebibliography}
\end{document}
